\section{Pathologies and Limits}


This section organizes the extent of validity of the framework in this paper and what lies beyond its scope.

"Pathologies" here does not mean medical diagnostic names.

This paper is not a theory for diagnosing psychiatric, neurological, or consciousness disorders.

What is treated here is structural deviation: the directions in which the log-referencing state can collapse when it departs from the consciousness window introduced in preceding sections.

The purpose of this section is therefore not to explain pathology but to clarify the limit conditions of conscious states.

\subsection{Pathology as Regime Deviation}


In the framework of this paper, stable conscious states are expressed as a window. Outside this window, consciousness does not merely "weaken." Collapse has direction:

\begin{itemize}
\item Direction of vanishing speed difference.
\item Direction of excessive speed difference.
\item Direction of too-low coupling viscosity.
\item Direction of too-high coupling viscosity.
\item Direction of too-weak meaning gradient.
\item Direction of too-strong meaning gradient.
\item Direction where live boundaries sediment excessively through measurement or report.
\end{itemize}

Each brings a different deformation of referencing state.

The crucial point is not linking these directly to disease names.

What this paper treats is not diagnostic classification but regime transition across time scales.

\subsection{Collapse Toward Synchronization}


The first deviation is the direction where speed difference becomes too small ($\Delta v \rightarrow 0$).

The boundary between fast and slow layers collapses into synchrony. Processing stabilizes and becomes efficient, but temporal thickness for referencing is lost.

In this direction, action automates, reactions become smooth, and explicit conscious referencing weakens.

This is not necessarily a bad state. Skilled movement, habituated action, and much reflexive processing are close to this direction.

But as one moves too far toward synchrony, reinterpretation of situations and re-referencing of logs becomes difficult. Consciousness does not so much stop as become unnecessary: processing runs, but the need for log referencing as opening diminishes.

\subsection{Dissociation by Excessive Difference}


The second deviation is the direction where speed difference becomes too large ($|\Delta v| \rightarrow \Delta v_{\max}$ or beyond).

The correspondence between fast and slow layers becomes unstable. Associations, predictions, images, and reactions run in the fast layer, but their stable sedimentation into the slow layer becomes difficult.

In this direction, the conscious state tends to fragment. Dreams, daydreaming, immersive states, dissociative states, and hallucinatory states can be described as part of this regime.

But care is needed here too. This paper does not equate these states. What they share is only the structure of increased speed difference and decreased stability of log referencing.

This section therefore indicates not diagnostic classification but the direction of referencing-structure breakdown.

\subsection{Fixation by Low Viscosity}


The third deviation is where coupling viscosity $\mu$ is too low ($\mu \rightarrow 0$).

Fast and slow layers adhere excessively. Speed difference does not slip sufficiently; processing tends to bind strongly to specific logs or meaning wells.

In this direction, states such as repetitive thought, hyperfocus, fixation, and compulsive referencing tend to arise.

The crucial point is that what is occurring here is not merely "strength of meaning."

The problem is that the coupling of meaning gradient and low viscosity causes processing to fall repeatedly into the same well.

Fixation is thus treated not as a content problem but as a viscosity problem of the referencing pathway.

Conscious content may feel vivid in this case, but vividness does not mean high degrees of freedom. Rather, fixation to meaning wells has strengthened and re-interpretability has decreased.

\subsection{Fragmentation by High Viscosity}


The fourth deviation is where coupling viscosity $\mu$ is too high ($\mu \rightarrow \infty$).

Fast and slow layers slip excessively against each other. Even though fast processing runs, it does not sufficiently sediment as a slow log.

Experience becomes difficult to retain; self-continuity and temporal coherence weaken.

In this direction, referencing fragmentation, sedimentation failure into memory, and reporting instability can occur.

This paper does not map this to a specific disease. What is stated here is only the structural tendency predicted when coupling viscosity is too high.

Processing exists, but does not stabilize as log referencing. This differs from collapse toward the synchrony side.

\subsection{Weak Meaning Gradient}


The fifth deviation is where the meaning gradient is too weak ($|\nabla\Phi| \rightarrow 0$).

Log referencing loses direction. Even though processing exists, which logs should be referenced and into which conceptual wells to fall becomes unclear.

In this direction, consciousness becomes diffuse, attention becomes scattered, and the outline of experience weakens.

But weak meaning gradient does not necessarily mean functional decline. Rest, diffuse thinking, free association, and creative wandering may utilize this region.

The problem arises when the meaning gradient is so weak that log referencing cannot form a stable terrain.

A weak meaning gradient can be soft exploration or loss of referencing — the difference is determined by the combination of speed difference, viscosity, delay, and gate state.

\subsection{Excessive Meaning Gradient}


The sixth deviation is where the meaning gradient is too strong ($|\nabla\Phi| \rightarrow \infty$).

Processing is strongly drawn into specific meaning wells. Log referencing becomes vivid, but degrees of freedom decrease.

In this direction, states such as emotionally strong memories, traumatic re-referencing, compulsive repetition, and excessive identification with narratives can be described.

The crucial point is that meaning gradient does not only strengthen consciousness.

Gradient directs referencing, but excessive gradient fixes referencing.

Strong meaning does not therefore necessarily mean deep understanding. An overly strong meaning field can rather take away the plasticity of understanding.

In this respect, $\mathcal{T}_{\Phi}$ is ambivalent: it is both the condition that opens log referencing and, in excess, the condition that fixes referencing.

\subsection{Measurement-Induced Collapse}


The seventh deviation is where live boundaries sediment excessively through measurement or report.

As stated in §13, measurement does not preserve qualia themselves.

Measurement sediments live boundary states into the slow layer and converts them into reportable residue.

When measurement demands are too strong, the conscious state tends to be treated not as a live opening but as an excessively stabilized log.

What remains is vivid: verbalized, quantified, made comparable. But that vividness is not live qualia themselves. It is the vividness of residue after qualia have sedimented.

This deviation is important in experimental design and introspection. The more one tries to observe, the more the object is deformed. The more one tries to explain, the more the live boundary is logged.

In the framework of this paper, measurement and introspection are therefore also treated as regime-changing operations that alter the conscious state.

\subsection{Mislocalization of Consciousness}


The eighth deviation is treating consciousness as identical to local substrates or sedimented residue.

As stated in §14, consciousness has neural substrates. But the conscious state itself does not localize to a single site.

What localization captures is local substrates, neural correlates, reportable logs, and post-measurement residue.

These are important. But identifying them with consciousness itself causes the live log-referencing regime to become invisible.

This error appears not only in discussions seeking the seat of consciousness but also in discussions seeking the self or qualia as specific objects within the brain.

This paper does not deny local substrates.

But local substrates can be necessary conditions for consciousness without being consciousness itself.

Consciousness is the state that opens when substrates enter a specific temporal, semantic, and non-closed relation.

\subsection{Limits of the Present Framework}


The descriptions so far are intended to organize deviations of conscious states structurally.

But this paper has clear limits.

First, this paper is not a medical diagnostic theory. The regimes described here do not directly correspond to psychiatric or neurological classifications.

Second, the measurable definition of $\mathcal{T}_{\Phi}$ is not yet sufficiently established. The relation $\mathcal{T}_{\Phi} \sim \kappa \frac{\Delta v}{\mu} |\nabla\Phi|$ is at present a structural expression, not a concrete measurable quantity.

Third, how to map the IFGT definitions of $\Phi$, $I$, $J$ onto neural activity, linguistic report, and subjective report remains incomplete.

Fourth, treating qualia as boundary phenomena does not explain all of subjective experience. It repositions the qualia problem without eliminating subjectivity as a whole.

Fifth, this paper does not provide a unified equation of consciousness. Attempting to fix consciousness in a closed equation loses the temporal non-closure that sustains it.

The theoretical scope of this paper lies therefore not in the complete description of consciousness itself but in the specification of the structural conditions under which consciousness can obtain.

\subsection{Why the Limits Matter}


These limits are both weaknesses of the paper and safety mechanisms of the theory.

Attempting to measure consciousness as a single variable immediately returns the theory to the generation model.

Treating qualia as stored objects loses their character as boundary phenomena.

Linking pathologies directly to diagnostic names loses the flexibility of structural regimes.

Placing consciousness at local sites makes the live log-referencing loop invisible.

This paper therefore deliberately leaves several points unclosed.

Not closing is not vagueness.

It is because the object this paper treats — consciousness — is itself a temporal structure maintained in non-closure.

If the theory rushes to complete closure, it loses the consciousness that is its object.

Making this constraint explicit is the purpose of this section.

\subsection{Summary}


Pathologies in this paper are not disease classifications but directions of deviation from the consciousness window.

\begin{itemize}
\item $\Delta v \rightarrow 0$: synchronization, automation, weakening of conscious referencing.
\item Excessive $|\Delta v|$: fragmentation, dream-like states, dissociation-like states, hallucinatory-like states.
\item $\mu \rightarrow 0$: fixation, repetition, fixing to meaning wells through low viscosity.
\item $\mu \rightarrow \infty$: sedimentation failure, fragmentation of referencing through high viscosity.
\item $|\nabla\Phi| \rightarrow 0$: disappearance of meaning gradient, weakening of outline.
\item Excessive $|\nabla\Phi|$: over-fixation to meaning wells, decrease of plasticity.
\item Excessive measurement demands: conversion of live qualia into sedimented residue.
\item Excessive localization: mistaking local substrates or residue for consciousness itself.
\end{itemize}

None of these determine the presence or absence of consciousness in binary terms.

They are structural classifications showing in which direction the conscious state deforms and under which conditions it leaves the stability window.
